\documentclass[10pt,a4paper]{amsart}
\usepackage[margin=19mm]{geometry}
\usepackage{amsmath,amssymb,mathtools}
\usepackage[scr=rsfs,cal=euler]{mathalfa}
\usepackage{xfrac}
\usepackage[unicode,hidelinks,bookmarks=false]{hyperref}
\newcommand{\ie}{i.e.\ }
\newcommand{\cf}{cf.\ }
\newcommand{\resp}{respectively\ }
\newcommand{\Subsection}{\subsection}
\providecommand{\nobreakdash}{\nobreak\hbox{-}}
\setlength{\emergencystretch}{2em}
\multlinegap0pt
\usepackage{amssymb}
\usepackage{enumerate}
\usepackage{bm}
\newtheorem{thm}{Theorem}
 \def\BB{{\mathbb B}}
 \def\CV{{\mathcal V}}
 \def\NN{{\mathbb N}}
 \def\RR{{\mathbb R}}
 \def\a{{\alpha}}
\def\f{\frac}
 \def\fD{{\mathfrak D}}
 \def\fE{{\mathfrak E}}
 \def\la{{\langle}}
 \def\ra{{\rangle}}
 \def\xb{{\bm x}}
 \def\yb{{\bm y}}
\title{Orthogonal polynomials which are eigenfunctions of a partial differential operator}
\author{Yuan Xu}
\date{Source: arXiv:2609.01751v1 (CC BY 4.0)}
\begin{document}
\maketitle

\noindent\emph{Source and licence.} Orthogonal polynomials which are eigenfunctions of a partial differential operator. Version arXiv:2609.01751v1 (CC BY 4.0), distributed by arXiv under the Creative Commons Attribution 4.0 International licence, http://creativecommons.org/licenses/by/4.0/, as recorded in the arXiv licence metadata for this exact version and shown on the abstract page https://arxiv.org/abs/2609.01751v1. Full licence notice: \texttt{licenses/arxiv-2609.01751-CC-BY-4.0.txt}.\par\smallskip

\noindent\textit{Editorial scope.} Counted unit: the article's thm thm:RB (source0.tex lines 2092--2106). The article's complete original proof of it is delivered with it (source0.tex lines 2108--2207, one \texttt{proof} environment of 100 lines). No mathematics is written, repaired or re-derived: the statement and the proof are the article's own lines, in the article's own notation, and no retained line is substituted or re-flowed.  Retained uncounted as the source-local context the proof needs: lines 479--481; lines 676--679; lines 1135--1137; lines 1139--1141; lines 2081--2084.  Nothing was omitted from any retained span, so the package carries no omission record.  No retained line contains a citation, so the wrapper carries no bibliography and no bibliography entry.  No retained line contains a figure environment, an image-inclusion command or drawing code, so no image is delivered and the package declares no asset dependency.  Editorial formatting only, in addition: an AMS \texttt{amsart} preamble that restates the article's own declarations and macros used by the retained lines, namely the environments \texttt{thm} on separate counters, together with the standard packages those lines need.  The retained environments print in this document as Theorem 1 in order of appearance, whereas the article prints the same items with the numbering of its own full document.  The complete native source of the article as distributed in the arXiv e-print for this exact version is bundled unchanged as \texttt{sources/209140/source0.tex}.
\begin{equation} \label{eq:Diff-ball}
\fD_\mu^{\BB} = \fD_\mu^{\BB,d} :=  \Delta_{\bm x} - \la \bm x, \nabla_{\bm x}\ra ^2 - (2\mu + d-1) \la \bm x, \nabla_{\bm x} \ra  
\end{equation}

\begin{equation}\label{eq:diff-eqnH}
   \fD^H u : = ( \Delta - 2 \la \xb, \nabla \ra )u = - 2 n u, 
    \quad u \in \CV_n\left(\RR^d, \bm W^H\right).  
\end{equation}

\begin{equation} \label{eq:diffHa}
  \rho(\yb)  \partial_{y_i} H = \partial_{y_i} \rho(\yb) \left( m  - \la \xb,\nabla_\xb\ra \right) H, \quad 1 \le i\le d_2,
\end{equation}

 \begin{equation} \label{eq:diffHb}
  \rho(\yb)  \la \yb, \nabla_{\yb}\ra  H = \la \yb, \nabla_{\yb}\ra \rho (\yb) \big( m  - \la \xb,\nabla_\xb\ra H \big). 
\end{equation}

 \begin{equation} \label{eq:OP_RB}
 \bm Q_{\bm j, \bm k, m}^n(\bm x,\bm y) = \BB_{\bm j, n-m}^{\mu+m+\frac{d_1}{2}}(\bm y) (1-\|\bm y\|^2)^{\frac{m}{2}}  
      \bm H_{\bm k, m}  \bigg(\frac{\bm x}{\sqrt{1-\|\bm y\|^2}} \bigg),
\end{equation}

\begin{thm}\label{thm:RB}
For $\mu > -\f12$, define the fourth-order differential operator 
\begin{align} \label{eq:RB}
  \fD_{\mu}^{\RR \rtimes \BB} u: = \,& -\frac14 (1-\|\yb\|^2) \Delta_{\xb}^2 + \la \xb,\nabla_\xb\ra \Delta_\xb - (\la \xb,\nabla_\xb\ra 
   + \la \yb,\nabla_\yb\ra) ^2  \\
   & + \Delta_\xb + \Delta_\yb % \left(\mu + \frac{d_1+d_2-1}{2}\right) \Delta_\xb    -
   +  (2\mu + d_1+d_2 -1) \left( \f12 \Delta_\xb- \la \xb,\nabla_\xb\ra - \la \yb,\nabla_\yb\ra \right). \notag
\end{align}
Then $\CV_n(\RR^{d_1}  \rtimes \BB^{d_2}, \bm W_{\mu}^{\RR \rtimes \BB})$ is an eigenspace 
of $ \fD_{\mu}^{\RR \rtimes \BB}$ for each $n \in \NN_0$. More precisely, 
\begin{align} \label{eq:RB-eigen}
  \fD_{\mu}^{\RR \rtimes \BB} u =  - n (n+ 2\mu + d_1+d_2 - 1) u, \qquad 
         \forall u \in \CV_n\big(\RR^{d_1}  \rtimes \BB^{d_2}, \bm W_{\mu}^{\RR \rtimes \BB}\big).
\end{align}
\end{thm}

\begin{proof}
We work with $u =\bm Q_{\bm j, \bm k, m}^n$ given in \eqref{eq:OP_RB} and write 
$u(\bm x,\bm y) = g(\bm y) H(\bm x,\bm y)$ with 
$$
g(\bm y) = \BB_{\bm j, n-m}^{\a+m}(\bm y) \quad \hbox{and} \quad
  H(\bm x, \bm y) =  (1-\|\bm y\|^2)^{\frac{m}{2}}  
      \bm H_{\bm k, m} \bigg(\frac{\bm x}{\sqrt{1-\|\bm y\|^2}}\bigg).
$$
where $\a = \mu+ \frac{d_1}2$.
With $\rho(\yb) = \sqrt{1-\|\bm y\|^2}$, it follows from \eqref{eq:diffHa} and \eqref{eq:diffHb} that
\begin{align} \label{eq:RB-diffH}
 \begin{split}
    (1-\|\yb\|^2) \partial_{y_i} H \, & = y_i (\la \xb, \nabla_\xb \ra-m) H, \quad 1 \le i \le d_2, \\
    (1-\|\yb\|^2) \la \yb, \nabla_\yb\ra H \,& = \|\yb\|^2 (\la \xb, \nabla_\xb \ra -m) H.
\end{split}
\end{align}
As a consequence of the two identities in \eqref{eq:RB-diffH}, we obtain readily that 
\begin{align*}
   (1-\|\yb\|^2)  \la \nabla_\yb g, \nabla_\yb H \ra  = \sum_{i=1}^{d_2} \partial_{y_i} g \cdot y_i  \,  (\la \xb, \nabla_\xb \ra  -m) H
    =  \, & \la \yb, \nabla_\yb \ra g (\la \xb, \nabla_\xb \ra -m) H, \\
(1-\|\yb\|^2) \la \yb, \nabla_\yb\ra g \cdot \la \yb \nabla_\yb \ra H 
     =  \| \yb\|^2 \, & \la \yb, \nabla_\yb\ra g  (\la \xb, \nabla_\xb \ra -m) H,
\end{align*}
so that their difference implies the identity
\begin{align} \label{eq:RB-diffH2}
 \la \nabla_\yb g, \nabla_\yb H \ra -\la \yb, \nabla_\yb\ra g \cdot \la \yb \nabla_\yb \ra H 
      =  \la \yb, \nabla_\yb \ra g  \big (\la \xb, \nabla_\xb \ra -m \big) H.  
\end{align}
Let $\fD_\a^{\BB, (\yb)}$ be the spectral operator \eqref{eq:Diff-ball} for the classical orthogonal polynomials on 
the unit ball in the variable $\yb$. Then, $\fD_{\a+ m}^{\BB, (\yb)} g - \fD_\a^{\BB, (\yb)} g= 2m  \la \yb, \nabla_\yb \ra g$,
so that a straightforward computation shows that the identity \eqref{eq:RB-diffH2} implies 
\begin{align*}
\fD_\a^{\BB, (\yb)} u -2 \la \yb,\nabla_{\yb} \ra \la \xb,\nabla_{\xb} \ra u 
     = \,& \fD_\mu^{\BB, (\yb)} g \cdot H + 2  \la \nabla_\yb g, \nabla_\yb H \ra 
         - 2 \la \yb, \nabla_\yb\ra g \cdot \la \yb \nabla_\yb \ra H \\
    & -2 \la \yb,\nabla_{\yb} \ra g \la \xb,\nabla_{\xb} \ra H + g \cdot \fE(H) \\
    = \, & \fD_{\a+ m}^{\BB, (\yb)}g\cdot H  + g \cdot \fE(H) ,
\end{align*}
where the operator $\fE$ is given by 
$$
  \fE = \Delta_\yb - \la \yb,\nabla_{\yb} \ra^2 -2 \la \yb,\nabla_{\yb} \ra \la \xb,\nabla_{\xb} \ra
  - (2\a + d_2 -1)  \la \yb,\nabla_{\yb} \ra.  
$$
Using $\fD_{\a+ m}^{\BB, (\yb)}g = - (n-m) (n+m + 2\a + d_2 -1)$, % = -n(n+2\a+d_2-1) + m (m+2\a+d_2 -1)$, 
we then obtain
\begin{align}\label{eq:Dmu_RB}
\fD_\a^{\BB, (\yb)} u -2 \la \yb,\nabla_{\yb} \ra \la \xb,\nabla_{\xb} \ra u 
     =\,&   -n(n+2\a+d_2 -1) u  \\
     &  + m (m+2\a+d_2 -1) u + g \cdot \fE(H), \notag
\end{align}

We now use \eqref{eq:RB-diffH2} to replace the derivatives with respect to $\yb$ in $\fE(H)$ by derivatives in
$\xb$. Applying the second identity in \eqref{eq:RB-diffH2} twice, a careful computation shows that 
\begin{align*}
 (1-\|\yb\|^2) \big(\Delta_\yb -   \la \yb,\nabla_{\yb} \ra^2\big)H =\, &
       (1-\|\yb\|^2) \la \yb,\nabla_{\yb} \ra ( \la \xb,\nabla_{\xb} \ra - m )H\\
  &   + 2 (1-\|\yb\|^2)  \la \yb,\nabla_{\yb} \ra H 
     + (d_2 - \|\yb\|^2) ( \la \xb,\nabla_{\xb} \ra -m )H \\
   = &\|\yb\|^2 
       ( \la \xb,\nabla_{\xb} \ra -m )^2 H 
     +  d_2 ( \la \xb,\nabla_{\xb} \ra -m )H.
\end{align*}
Substiting this identity to $(1-\|\yb\|^2) \fE(H)$ leads to, after further computation, that 
\begin{align} \label{eq:EH-R}
 (1-\|\yb\|^2) \fE(H) = \, & \|\yb\|^2 \left(-\la \xb,\nabla_{\xb} \ra^2 - (2\a -1) \la \xb,\nabla_{\xb} \ra+ m(m+2 \a -1) \right)H \\
     & + d_2 (1-\|\yb\|^2) (\la \xb, \nabla \xb\ra -m)H. \notag
\end{align}
By the definition of $H$ and the spectral equation \eqref{eq:diff-eqnH}, we obtain
\begin{equation} \label{eq:diffH-RB}
    (1-\|\yb\|^2) \Delta_\xb H - 2 \la \xb, \nabla_\xb \ra H = - 2 m H,
\end{equation}
which can be rewritten as 
$$
   m H = \la \xb, \nabla_\xb \ra H - \frac12 (1-\|\yb\|^2) \Delta_{\xb} H.
$$
Using this relation twice, we deduce 
\begin{align} \label{eq:mm-H}
   & m(m+ 2 \a -1) H  = \la \xb, \nabla_\xb \ra^2 + (2\a-1)  \la \xb, \nabla_\xb \ra \\
    &\,  - \frac12 (1-\|\yb\|^2) \left[  \la \xb, \nabla_\xb \ra \Delta_\xb + \Delta_\xb \la \xb, \nabla_\xb\ra 
      - \frac 12(1-\|\yb\|^2) \Delta_\xb^2 + (2\a-1) \Delta_\xb \right], \notag
\end{align} 
which implies an expression of $\fE(H)$ given by 
\begin{align*}
\fE(H) = - \frac{\|\yb\|^2}2  \left[  \la \xb, \nabla_\xb\ra \Delta_\xb + 
    \Delta_\xb  \la \xb, \nabla_\xb\ra - \frac 12(1-\|\yb\|^2) \Delta_\xb ^2 
       + (2\a-1) \Delta_\xb  \right] H  \\
        + \f{d_2}{2} (1-\|\yb\|^2) \Delta_x  H.
\end{align*}
Using this identity and \eqref{eq:mm-H} with $\a$ replaced by $\a+ \frac{d_2}2$, we obtain from \eqref{eq:Dmu_RB} that 
\begin{align*}
& \fD_\a^{\BB, (\yb)} u  -2 \la \yb,\nabla_{\yb} \ra \la \xb,\nabla_{\xb} \ra u - \la \xb,\nabla_\xb \ra^2  u -(2\a-1) \la \xb,\nabla_\xb \ra u\\
&\qquad +  \frac{1}2  \left[  \la \xb, \nabla_\xb\ra \Delta_\xb  + \Delta_\xb  \la \xb, \nabla_\xb\ra - \frac 12(1-\|\yb\|^2) \Delta_\xb ^2 
       + (2\a-1) \Delta_\xb  \right] u \\
      &\quad  =   -n(n+2\a+d_2 -1) u= -n (n+ 2\mu + d_1+d_2-1) u,
\end{align*}
after simplification via $\|\yb\|^2 + (1-\|\yb\|^2) =1$. This is \eqref{eq:RB}, and its left-hand side can be 
simplified using the expression of $\fD_\a^{\BB, (\yb)}$ in \eqref{eq:Diff-ball}, as well as $\a = \mu + \frac{d_1}2$ and 
$\Delta_\xb \la \xb, \nabla_\xb\ra = \la \xb, \nabla_\xb\ra \Delta_\xb + 2 \Delta_\xb$, 
to the expression $\fD_\mu^{\RR \rtimes \BB}$ in \eqref{eq:RB}. 
\end{proof} 

\end{document}
